\documentclass[11pt]{article}
\usepackage[margin=1in]{geometry}
\usepackage[T1]{fontenc}
\usepackage{lmodern,amsmath,amssymb,booktabs,graphicx,enumitem}
\usepackage[round,authoryear]{natbib}
\usepackage[colorlinks=true,linkcolor=blue,urlcolor=blue,citecolor=blue]{hyperref}
\setlength{\parindent}{0pt}
\setlength{\parskip}{5pt}
\setlength{\emergencystretch}{2em}
\title{RoCE-K: residual compatibility experiments}
\author{}
\date{1 October 2026}
\begin{document}
\maketitle
\textbf{Scope.} These are oracle residual-score experiments and a mathematical
implementation audit. They do not validate fitted high-dimensional RoCE,
cross-fitting with overlapping training sets, or a full growing-source theorem.
The reported design has 114 settings and 138,000 setting-specific repetitions.
The original six extra-pilot settings are archived separately and excluded from
the final panel. No historical results or production estimators were changed.

\section{Research question and experimental decisions}

We test whether source-residual compatibility can jointly protect against weak
bias and increasing source counts. The target decomposition is
\[
 \tau=d+r_1-r_0,\qquad
 d=E_t(m^1-m^0),\quad r_a=\mu_t^a-E_tm^a.
\]
The target score is $(m^1-m^0,\psi_t^1-m^1,\psi_t^0-m^0)$, so the shared
target component is included only once. The source module supplies a set for
each $r_a$. Intersect those constraints with the joint target region, then
project onto $(1,1,-1)$. Source valid sets may differ by arm.

Searching under uncertain validity is motivated by \citet{guo2023}; the
multiscale frequency shape is motivated by \citet{wang2026}. Our count-search
baseline is not a reproduction of Guo's searching-and-sampling algorithm.
Our fixed site means, heterogeneous variances and target localization also
differ from the latter paper's Gaussian mixture experiment. Its optimality
and length theorems are not asserted for this prototype.

The initial 90 settings use 1,000 repetitions each. The main Gaussian panel
has $K=8,32,128,512,2048$, two target shared-noise levels, and four bias
patterns. Sources in a fixed quarter of the sites have treated-arm shifts
equal to $0$, $2\sigma K^{-1/3}$, $2\sigma K^{-1/4}$ or $4\sigma$, where
$\sigma$ is a source residual standard error. Additional panels cover
different valid arm sets, equal-arm shifts, one extreme source, heterogeneous
variances, weaker valid-count guarantees, pilot variance estimation and binary
patient scores. Main experiments declare at least $75\%$ valid sources in
each arm; their identities are not supplied to the procedures.

After the main Gaussian results, 18 confirmation settings were declared with
2,000 new-seed repetitions each. They compare an error-budgeted intersection
of Fourier and dispersion certificates. Six further settings, also with
2,000 new-seed repetitions, have actual and declared validity fractions of
$1/2$ or $1/4$. This is a boundary study under valid assumptions, not a
violation of the validity-count guarantee.

\section{Inference and numerical safeguards}

For independent Gaussian residual means $X_j\sim N(\theta_j,v_j)$,
the corrected empirical characteristic function is
\[
 \widehat F(t)=K^{-1}\sum_j\exp(itX_j+t^2v_j/2).
\]
With $\rho=g/K$, candidate $r$ is retained when
$|\widehat F(t)-\rho e^{itr}|\le1-\rho+\Delta(t)$ at every frequency.
The grid is $t_\ell=\sqrt\ell/(2\sqrt{v_{\max}})$ for
$\ell=1,\ldots,\lceil\log(eK)\rceil$. It is specified before evaluation;
no frequency is selected for favorable coverage.

Let $m$ be the number of frequencies, $x=\log(4m/\alpha_s)$,
$M_j=\exp(t^2v_j/2)$, $V=\sum_j(M_j^2-1)$ and $B=\max_j(M_j+1)$.
Two valid Gaussian noise allowances are
\begin{align*}
 \Delta_H(t)&=\frac2K\sqrt{x\sum_jM_j^2},\\
 \Delta_B(t)&=\frac{\sqrt2}{K}
 \left\{\frac{Bx}{3}+\sqrt{2Vx+(Bx/3)^2}\right\}.
\end{align*}
The first uses Hoeffding. For the second, each real or imaginary component
has variance at most $M_j^2-1$, and its centered absolute value is at most
$M_j+1$. Apply Bernstein to both tails of both components and union across
frequencies. The smaller of these deterministic bounds can be selected before
evaluation without another error allocation. This tightening is not an
empirically tuned critical value.

All Fourier disks are inverted analytically into periodic intervals and
intersected as interval sets. No candidate grid is used to measure coverage.
The target ellipsoid is projected over every residual-interval rectangle.
The unconstrained optimum and four analytically optimized edges give its
endpoints; disconnected source sets are retained through projection.
Numerical endpoints receive small outward padding. Independent constrained
optimization and dense disk-membership checks verify these calculations.

The target region uses failure allowance $.025$, and each source arm receives
$.0125$. In the combined method, each arm splits that allowance equally
between Fourier and dispersion. This is an intersection of simultaneously
protected constraints, not selection of the shortest unadjusted interval.
The dispersion comparator uses noncentral chi-square inversion only for
homogeneous exact Gaussian noise. Its heterogeneous version uses a moment
bound. Empty feasible sets return the prespecified target interval with a
flag; no repetition is dropped. Midpoints of the resulting intervals are
reported as exploratory point estimates.

\section{Results in the Gaussian oracle experiment}

The following table uses independent confirmation data: $K=2048$, no shared
target variance floor, and 2,000 repetitions per setting. Length is relative
to the ordinary $95\%$ target-only interval. The valid-set oracle has extra
information and is explicitly separated from implementable comparators.

\begin{center}\small
\input{confirmation_table}
\end{center}

At weak shifts, dispersion is shorter than Fourier in this known-variance
model. At strong shifts, Fourier is shorter. The combined certificates retain
most of the weak-bias gain and improve the strong-bias result, in part because
different certificates can be more useful in different arms. Most protected
intervals overcover relative to $95\%$; the experiment does not establish
optimal calibration or oracle efficiency. Source counts this large are a
growth stress test, not a claim about typical federated study sizes.

When shared target noise is present, it can mask the weak bias in a naive
TATE interval and limit the gain from sources. This differs from correctness
of the source residual module itself. Small-$K$ protected intervals can be
longer than ordinary target-only intervals, because their joint-region and
source-uncertainty protection is costly.

\paragraph{Point estimation and fallback require a separate repair.}
For the combined interval at $K=2048$ without a target variance floor,
the all-valid confirmation RMSE is $.00663$, versus $.00100$ conditional
on no fallback. The $.8\%$ fallback cases account for $97.7\%$ of squared
error. At weak shifts the corresponding share is $92.0\%$, and at strong
shifts $93.8\%$. Therefore a shorter interval does not establish a better
point-estimation rule. A diagnostic projection of naive pooling into the
interval helps some weak-bias cases but worsens strong-bias RMSE; it is not
adopted as a universal fix.

There is also an asymptotic issue. In the exact Gaussian zero-floor setting,
if source residual sets shrink to their true pair while the target ellipsoid
keeps a fixed error allowance, a nonvanishing probability of target-region
incompatibility can remain. Returning a target-width interval on that event
can prevent the desired expected-length gain in $K$. Shrinking the target
failure allocation or constructing a suitable shorter fallback are research
options requiring their own validation. The present fallback is retained
in every reported result; neither failures nor their MSE contributions are
removed.

\begin{figure}[p]
\centering
\includegraphics[width=\linewidth]{gaussian_main.pdf}
\caption{Initial Gaussian panel. Shading gives pointwise Monte Carlo Wilson
intervals for coverage. The two shared-noise levels use the same underlying
random draws. Each setting has 1,000 repetitions and $n=1000$ per site.}
\end{figure}
\begin{figure}[p]
\centering
\includegraphics[width=\linewidth]{independent_confirmation.pdf}
\caption{Independent confirmation, 2,000 repetitions per setting. Combined
certificates use a prespecified split of the source error allowance.}
\end{figure}

\section{Estimated variance and non-Gaussian scores}

\paragraph{Patient budget.}
Every reported setting has 1,000 patients per site. The source variance panel
uses 250 independent pilot patients and 750 evaluation patients; target
summaries use 1,000 patients and known target covariance. The initial version
of those six cells used 1,000 evaluation observations plus 250 additional
pilot observations. They are retained as auxiliary-information diagnostics
but replaced in the reported panel. Source variances and local shifts were
recomputed for the corrected evaluation size.

Gaussian pilot covariance is generated with the correct two-arm Wishart
dependence. Simultaneous chi-square variance intervals receive total failure
allowance $.005$. For certified lower and upper mean variances $l_j,u_j$,
deconvolution uses $l_j$ and adds the attenuation bound
$K^{-1}\sum_j\{1-\exp[-t^2(u_j-l_j)/2]\}$. The two source error allowances
are then $.01$ each. Estimated-variance plug-in is separately labeled as a
diagnostic without this certificate.

\paragraph{Binary observed-data model.}
Each site has $X\in\{-1,1\}$ with equal probabilities and randomized
treatment probability $1/2$. Common predictions are
$m^1(X)=.55+cX/2$ and $m^0(X)=.45-cX/2$, with $c=0$ or $.3$.
Target outcomes follow those Bernoulli means, so TATE is exactly $.10$.
Selected sources have additive arm shifts. The oracle merged weights are
two. Multinomial counts of the eight $(X,A,Y)$ states generate the exact
patient-score means and unbiased variances, without materializing every row.

For a patient score of range $B_j$, variance $s_j^2$, and sample size $n_j$,
a telescoping characteristic-function comparison with centered Gaussian
summands gives the deconvolved error bound
\[
 e^{t^2v_j/2}\frac{|t|^3}{6n_j^2}
 \{B_js_j^2+2\sqrt{2/\pi}\,s_j^3\},\qquad v_j=s_j^2/n_j.
\]
The bounded-score method adds the average of these bounds over sources. Its
complex-exponential variance bound is $M_j^2\min(1,t^2v_j)$ and its centered
bound is $2M_j$. The target uses a second-moment ellipsoid, rather than an
unproved exact Gaussian pivot. This is deliberately a conservative reference.
The Gaussian-approximation versions are also shown, but good empirical
coverage does not certify them uniformly for growing $K$.

The binary experiment exposes a practical obstacle: with a nonzero target
variance floor, the fully protected target moment region can make the final
interval substantially wider than nominal target-only inference. Tightening
the target calibration is therefore as important as tightening the source
module. Similarly, valid pilot variance bounds carry a visible length cost.

\begin{figure}[p]
\centering
\includegraphics[width=\linewidth]{binary_scores.pdf}
\caption{Binary patient-score experiment at weak shifts, 1,000 repetitions
per setting. Fourier (Bernstein) here uses a Gaussian approximation; bounded
protection adds a distributional error bound and a target moment region.}
\end{figure}

\section{Boundary, reliability and next decision}

The actual-minority panel puts either $1/2$ or $1/4$ of treated-source centers
at the target residual, and the remainder one source standard error away.
Both centers have at least the declared mass. They therefore both satisfy
the population source compatibility constraints at every frequency. Target
information is required to distinguish them. Observed borrowing gains are
weaker with the lower valid fraction; the experiment does not claim the same
uniform shrinking rate as under strict majority validity.

Nine automated tests cover analytic arc inversion, independent numerical
projection, singular prediction components, exact distribution comparisons,
calibration simulations, minority ambiguity and patient budgets. Every saved
interval, coverage value, length and RMSE was independently recomputed during
the result audit. Paired length differences and Monte Carlo standard errors
are exported. No complete-method theorem is inferred from these checks.

The next design decision is to retain residual compatibility as the main
framework and study a prespecified combination of Fourier and dispersion
certificates. Priorities are useful target and source calibration under
estimated distributions, valid nuisance-bias allowances, and a joint
$(n,K)$ length theorem. Full high-dimensional nuisance fitting should enter
after those components remain useful under their uncertainty costs.

All raw records, frozen source versions, protocols, the variance-budget
correction and audit receipts are stored under the experiment's scratch
directory. Older manuscript designs and results remain available.

\bibliographystyle{plainnat}
\bibliography{references}
\end{document}
